\section{引言：为什么软件开发 Agent 正在崛起}
\subsection{行业背景与平台}
Graham Neubig 先以 CMU 教授和 All Hands AI 的首席科学家的身份开场，讲述他在软件开发与人工智能之间的双重积累。All Hands AI 正在维护开源的 OpenHands（前身 OpenDev）框架，目标是把“写代码”的能力包装成能够靠代理执行的线性流程，让开发者能专注于需求与评审，而把重复的测试、搜索与补丁交给 Agent。

他指出开放平台让社区可以共建调试工具与示例仓库，从实际 issue 中提取输入、在公开仓库中跟踪 Agent 操作，再把改进反馈给模型。这种闭环让研究能够直接走向工程实践，而不是停留在实验室的基准上。

\subsection{赋能普罗大众}
他特别提到 2011 年 Mark Andreessen 的文章 ``Why Software Is Eating the World''，指出各行各业正在被软件重新定义：从娱乐、农业、国防到医疗服务，软件已经逐步取代了物理流程，未来软件负责的领域只会越来越多。这也说明了为什么需要让更多人拥有快速构建软件的能力，而不是让软件开发继续被少数专业工程师垄断。

在介绍动机时，Graham 用近年来多个由软件系统驱动的科研突破作类比，强调当代科学进展正在与代码深度捆绑。对他而言，Agent 的价值在于能把工程能力下放到更广泛的研究者、产品经理与跨学科团队，让“软件创造”成为普遍能力而非稀缺资源。

\begin{knowledgebox}{Mark Andreessen 的预言}
2011 年的 ``Why Software Is Eating the World'' 认为：软件将成为所有产业服务的新基础设施，而自动化构建软件的能力则是决定竞争力的关键。Graham 整个演讲试图把这句话落地到 Agent 工程中。
\end{knowledgebox}

\begin{importantbox}{Agent 的价值主张}
软件开发 Agent 的核心承诺是：将具备代码理解、变更、验证与反馈能力的多模态体系，变成能够像人类开发者一样持续迭代产品的执行者。
\end{importantbox}

\subsection{本章小结}
\begin{itemize}
    \item 软件已经主导大部分行业，构建能力的稀缺性推动了 Agent 的实际需求。
    \item Graham 通过诺贝尔奖与软件领域的对照提醒我们：软件创造已经成为人类科学的核心技能之一。
\end{itemize}

